\kbvcontinued ಆಗುವುದು. ಹೀಗಾದಲ್ಲಿ ಸಂಧಿಪದವು ಅರ್ಥರಹಿತವಾಗುವುದು. ಇಂತಹ ಸಂದರ್ಭದಲ್ಲಿ ಎರಡೂ ಸ್ವರಗಳ ಮಧ್ಯೆ ಹೊಸದೊಂದು ಅಕ್ಷರ ಸೇರಿ ಸಂಧಿ- ಕಾರ್ಯವಾಗುವುದು. ಮೇಲಿನ ಉದಾಹರಣೆಯಲ್ಲಿ “ಯ್” ಎಂಬ ಅಕ್ಷರವು ಹೊಸದಾಗಿ ಬಂದು ಸೇರಿರುವುದನ್ನು ನೋಡಬಹುದು. ಹೀಗೆ\par

\indent ಸಂಧಿಕಾರ್ಯ ನಡೆಯುವಾಗ ಒಂದು ಅಕ್ಷರವು ಹೊಸದಾಗಿ ಬಂದು ಸೇರುವುದಕ್ಕೆ “ಆಗಮ ಸಂಧಿ” ಎಂದು ಹೆಸರು.\par

\begin{kbvnote}
\noindent{\kbvsecfont ವಿವರಣೆ :}\hspace{0.5em}ಮೇಲಿನ ಉದಾಹರಣೆಯಲ್ಲಿ ಯ-ಕಾರವು ಹೊಸದಾಗಿ ಬಂದು ಸೇರಿದೆ; ಆದುದರಿಂದ ಇದನ್ನು ಯ-ಕಾರಾಗಮ ಸಂಧಿ ಎಂದು ಹೇಳುತ್ತೇವೆ.
\end{kbvnote}

\indent ಬಾಣ + ಅನ್ನು = ಬಾಣವನ್ನು ; ಇಲ್ಲಿ “ವ” ಕಾರವು ಹೊಸದಾಗಿ ಬಂದು ಸೇರಿ ಸಂಧಿಕಾರ್ಯವಾಗಿದೆ. ಇದು ವ-ಕಾರಾಗಮ ಸಂಧಿ.\par

\vspace{3pt plus 2pt}\noindent\hspace*{1.2em}\kbvlabel{ಉದಾ}\par\vspace{1pt}

\begin{kbvrule}
\noindent\kbvbullet\hspace{0.5em}ಯ-ಕಾರ ಆಗಮ ಸಂಧಿ \ensuremath{\square} ವ-ಕಾರ ಆಗಮ ಸಂಧಿ
\end{kbvrule}

\begin{kbvderiv}
\vspace{4pt plus 3pt minus 1pt}
\noindent{\renewcommand{\arraystretch}{1.24}\setlength{\tabcolsep}{5pt}
\begin{tabularx}{\textwidth}{X X}
ಹಣೆ + ಅಲ್ಲಿ = ಹಣೆಯಲ್ಲಿ & ಮಗು + ಇಗೆ = ಮಗುವಿಗೆ \\
ಕೊಠಡಿ + ಇಂದ = ಕೊಠಡಿಯಿಂದ & ಗೋ + ಇಂದ = ಗೋವಿಂದ \\
ಮೈ + ಅನ್ನು = ಮೈಯನ್ನು & ಹೆಣ + ಅನ್ನು = ಹೆಣವನ್ನು \\
ದಾರಿ + ಅನ್ನು = ದಾರಿಯನ್ನು & ಹೂ + ಇನ = ಹೂವಿನ \\
\end{tabularx}}
\vspace{4pt plus 3pt minus 1pt}
\end{kbvderiv}

\kbvsechead{(೩)\hspace{0.45em}ಆದೇಶ ಸಂಧಿ :}

\indent ನಾವೆಲ್ಲರೂ ಜೊತೆಗೂಡಿ ಹೊರಡೋಣ ಎಂಬ ವಾಕ್ಯದಲ್ಲಿ ಜೊತೆಗೂಡಿ ಎಂಬ ಸಂಧಿ ಪದವನ್ನು ಬಿಡಿಸಿ ನೋಡೋಣ.\par

\begin{kbvderiv}
\vspace{4pt plus 3pt minus 1pt}
\begin{center}\begin{adjustbox}{max width=\textwidth}
{\renewcommand{\arraystretch}{1.22}\setlength{\tabcolsep}{4pt}
\begin{tabular}{l c l}
ಜೊತೆ + ಕೂಡಿ & {\color{kbvgrey}$=$} & ಜೊ + (ತ್ + ಎ) + (ಕ್ + ಊ) + ಡಿ \\
ಜೊತೆಗೂಡಿ & {\color{kbvgrey}$=$} & ಜೊ + (ತ್ + ಎ + ಗ್ + ಊ) + ಡಿ \\
\end{tabular}}
\end{adjustbox}\end{center}
\vspace{4pt plus 3pt minus 1pt}
\end{kbvderiv}

\indent ಇಲ್ಲಿ ಸ್ವರದ ಮುಂದೆ ವ್ಯಂಜನ ಬಂದು ಸಂಧಿಕಾರ್ಯವಾಗುವುದನ್ನು ಕಾಣಬಹುದು. ಇಲ್ಲಿ ಉತ್ತರಾರ್ಥ ಭಾಗದ ಮೊದಲಿಗಿರುವ “ಕ್” ವ್ಯಂಜನಕ್ಕೆ ಬದಲಾಗಿ ಸಂಧಿ ಪದದಲ್ಲಿ “ಗ್” ವ್ಯಂಜನ ಸೇರಿ ಸಂಧಿಕಾರ್ಯವಾಗಿರುವುದು.\par
